%!TEX root = ../dokumentation.tex

\newcommand{\pdftitel}{Comparison of a Classical RAG Pipeline and an Agent-Based Approach for Intelligent Knowledge Retrieval Across Enterprise Documentation Sources}
\newcommand{\autor}{Liam Thomson}
\newcommand{\arbeit}{Bachelor Thesis}

% Seitengröße
\usepackage{fullpage}

% Zeilenumbruch und mehr
\usepackage[activate]{microtype}

% Zeichencodierung
\usepackage[utf8]{inputenc}
\usepackage[T1]{fontenc}

% Zeilenabstand
\usepackage[onehalfspacing]{setspace}

% Index-Erstellung
\usepackage{makeidx}

% Language
\usepackage[english]{babel}

% Quotation marks
\usepackage[babel]{csquotes}

% Spezielle Tabellenform für Deckblatt
\usepackage{longtable}
\setlength{\tabcolsep}{10pt}
\renewcommand{\arraystretch}{1.5}

% Grafiken
\usepackage{graphicx}

% Farben
\usepackage[dvipsnames]{xcolor}
\definecolor{LinkColor}{rgb}{0,0,0.2}
\definecolor{ListingBackground}{rgb}{0.92,0.92,0.92}

% Titel, Autor und Datum
\title{\titel}
\author{\autor}
\date{\datum}

% PDF Einstellungen
\usepackage[%
	pdftitle={\pdftitel},
	pdfauthor={\autor},
	pdfsubject={\arbeit},
	pdfcreator={pdflatex, LaTeX with KOMA-Script},
	pdfpagemode=UseOutlines,
	pdfdisplaydoctitle=true,
	pdflang=en
]{hyperref}

\hypersetup{%
	colorlinks=true,
	linkcolor=LinkColor,
	citecolor=LinkColor,
	filecolor=LinkColor,
	menucolor=LinkColor,
	urlcolor=LinkColor,
	bookmarksnumbered=true
}

% Literaturverweise nach Harvard
\usepackage[dcucite]{harvard}

% Schriftart
\usepackage{palatino}

% Hurenkinder und Schusterjungen verhindern
\clubpenalty=10000
\widowpenalty=10000
\displaywidowpenalty=10000

% Quellcode
\usepackage{listings}
\lstloadlanguages{Python}
\lstset{%
	language=Python,
	stepnumber=1,
	numbersep=5pt,
	numberstyle=\tiny,
	breaklines=true,
	breakautoindent=true,
	postbreak=\space,
	tabsize=2,
	basicstyle=\ttfamily\footnotesize,
	showspaces=false,
	showstringspaces=false,
	extendedchars=true,
	captionpos=b,
	backgroundcolor=\color{ListingBackground}
}

% Glossar
\usepackage[
	nonumberlist,
	toc,
]{glossaries}

% Akronyme
\usepackage[printonlyused,footnote]{acronym}

% Fußnoten
\usepackage[perpage, hang, multiple, stable]{footmisc}

% Bildpfad
\graphicspath{{images/}}

\usepackage{bookmark}

\usepackage{float}

\usepackage{array}

\usepackage{booktabs}

\usepackage[right]{eurosym}

\usepackage{wrapfig}

\usepackage{comment}
\specialcomment{k}{\begingroup\color{red}}{\endgroup}
%\excludecomment{k}
